

\begin{itemize}
\item \textbf{(A) Shape dimension.}
  The number of \(\bm{r}\)-dependent patterns (columns of \(\Phi\)) you need after removing rigid modes to span the Saint–Venant residual over the section. 
  Formally, it’s the separation rank of the residual \(R(x,\bm{r})\in \mathcal F_x\otimes(\mathcal V_r/\mathscr R)\): the smallest \(N\) such that
  \(R(x,\bm{r})=\sum_{j}^N \alpha_j(\reflenx)\,\bm{\Phi}_j(\bm{r})\) with \(\bm{\Phi}_j\) independent of \(\reflenx\).
  For a generic section: \textbf{8} (or \textbf{9} if you keep the uniform axial \(\reflenx^2\) term as a ``shape").

  Formally, \(\ShapeRank := \operatorname{rank}_{\text{sep}}(R)\); ie, the minimal number of \(\reflenx\)-independent warping shapes after rigid‑mode quotient.

  From the ``matrix–vector" viewpoint \(\bm\Phi(\bm{r})\,\bm\alpha(\reflenx)\), consider \(\bm\Phi(\bm{r})\in\mathbb R^{(3|\Section|)\times \ShapeRank}\): columns are fixed warping shapes \(\bm{\Phi}_j(\bm{r})\) after quotienting rigid modes.
  The shape rank is the number of columns \(\ShapeRank\) in a minimal such \(\bm\Phi\).

\item \textbf{(B) Field Rank.}
  The number of independent scalar functions of \(x\) required after assocciating as many amplitude channels as possible to the 1D beam measures \((\epsilon,\bm\gamma,\bm\kappa)\). 
  Concretely, if some shape’s required amplitude \(\alpha_j(\reflenx)\) has exactly the same \(x\)-law as a component (or fixed linear combination) of \(\epsilon,\gamma,\kappa\), you can set \(\alpha_j(\reflenx)\) equal to that measure—so it does not count as a new unknown function. 
  What’s left are the unavoidable amplitude functions. For a generic section: \textbf{5}.

  Formally, \(\FieldRank := \dim\big(\mathcal G/\mathcal E\big)\), where \(\mathcal G\) is the space spanned by all required amplitude functions \(\alpha_j(\reflenx)\) and \(\mathcal E\) is the subspace generated by fixed linear combinations of (\(\epsilon,\gamma,\kappa\)). 
  Equivalently, it is the count of amplitude channels whose \(x\)-laws cannot be written as such combinations for all loads.
  
From the ``matrix–vector" viewpoint \(\bm\Phi(\bm{r})\,\bm\alpha(\reflenx)\),  \(\bm\alpha(\reflenx)\in \mathbb R^\ShapeRank\): entries are amplitude channels \(\alpha_j(\reflenx)\).
  You now impose identifications of the form
  \[
  \alpha_j(\reflenx)=\mathcal L_j\big[\epsilon(\reflenx),\bm\gamma(\reflenx),\bm\kappa(\reflenx)\big]
  \]
  with \(\mathcal L_j\) a fixed, geometry‑dependent linear mapping that must work for all load parameters. Every channel that admits such an identification is not an independent unknown function. 
  The \emph{free‑amplitude rank} is the number of entries of \(\bm\alpha(\reflenx)\) that cannot be written in this way.

\end{itemize}

These numbers differ because many shapes can share the same driving signal in \(x\). 
The shapes control where the field lives across the section; the amplitudes control how it varies along \(x\). 
You may need more distinct shapes than you need distinct \(x\)-signals to drive them.

We make this precise and get a clean lower bound by treating the SV field as a separated function of the axial variable \(\reflenx\) and the in-plane variable \(\bm{r}\), and then quotienting out what the kinematic fields \((\bm x,\bm\theta)\) can already generate.

\paragraph{Rigid Quotient and the Shape Rank}


\begin{proposition}[Shape Rank]\label{prop:min8}
For a generic cross–section (no symmetry assumptions, arbitrary \(\CentroidLocation\)), the full Saint–Venant family
\((\IesanLoadM,a_\varepsilon,a_\vartheta,\bm{b}_{\Section})\) can be represented minimally with \(\ShapeRank = 8\) \(\reflenx\)-independent warping shapes of \(\bm{r}\).
\end{proposition}

\textbf{Proof}

\subparagraph{Proof Summary}
An outline is:
\begin{enumerate}
    \item \textbf{Enumerate monomials in \(\reflenx\):} \( {1,x,\reflenx^2,\reflenx^3}\).
  \item \textbf{Assemble SV coefficients:} For each monomial and each parameter direction \((a_y,a_z,a_\varepsilon,a_\vartheta,b_y,b_z)\), extract the corresponding \(r\)-shape \( \mathbf{S}_{m,p}(\bm{r})\).
  \item \textbf{Project out kinematics:} Build the projection onto \(\mathscr{R}(\Section)\) by least-squares (or exact algebra) using the rigid basis \(\{\PlaneBasis_y,\,\PlaneBasis_z,\,\FrameBasis\times\bm{r},\,\FrameBasis\otimes\bm{r},\,\FrameBasis\}\); subtract to obtain residual shapes \( \mathbf{S}^{\perp}_{m,p}\).
  \item \textbf{Block the \(\reflenx\)–dependence:} Stack the residuals by monomial: \(\mathbf{R}(x,\bm{r})=\sum_m f_m(\reflenx)\mathbf{R}_m(\bm{r})\) where \(\mathbf{R}_m\) is the matrix (over parameter directions) of shapes for that monomial.
  \item \textbf{Compute separation rank:}
  Form the span of the columns of all \(\mathbf{R}_m\) in \( \mathcal{V}_r/\mathscr{R}(\Section)\). 
  The dimension of this span is \(\ShapeRank\). 
  % Numerically, discretize \(\Section\), assemble a Gram matrix of these residual shapes (after removing rigid components), and take its rank (or SVD with tolerance).
\end{enumerate}

Write the full Saint–Venant displacement \eqref{eq:iesan-general-solution} as a sum of ``monomial in \(\reflenx\)" \(\times\) ``shape in \(\bm{r}\)":

$$
\hat{\bm u}(\reflenx,\bm{r})=\sum_{m\in\{0,1,2,3\}} f_m(\reflenx)\,\mathbf{S}_m(\bm{r}),
$$
where the \(\mathbf{S}_m(\bm{r})\) collect all \(\bm{r}\)-dependent patterns multiplying \(1,x,\reflenx^2,\reflenx^3\) in the full SV expression. 
Now split the kinematic image generated by
$$
\mathscr{R}(\Section)=\left\{\,\bm u(\reflenx)+\bm\theta(\reflenx)\times\bm{r} \;\middle|\; \bm u,\bm\theta \text{ arbitrary smooth in } \reflenx \,\right\}.
$$
\(\mathscr{R}(\Section)\) is spanned (pointwise in \(\reflenx\)) by rigid \(\bm{r}\)-patterns which preserve plane section:
\begin{equation}
\mathscr{R}(\Section) = \operatorname{span}
\underbrace{\{ \PlaneBasis_y,\,\PlaneBasis_z,\,\FrameBasis\times\bm{r}\}}_{\text{Transverse}}
\oplus 
\underbrace{\{\FrameBasis\otimes\bm{r},\,\FrameBasis\}}_{\text{Axial}}.
\label{eq:rigid-span}
\end{equation}
Anything in this span can be carried by the alignement fields $(\bm{u}, \bm{\theta})$ and is therefore removed before counting warping shapes. Everything outside this span requires explicit warping shapes and amplitudes.
Project the SV field onto \(\mathscr{R}(\Section)\) (choose \(\bm u,\bm\theta\) to cancel those parts), and call the residual
$$
\mathbf{R}(\reflenx,\bm{r}) \in \mathcal{F}_x \otimes \big(\mathcal{V}_r / \mathscr{R}(\Section)\big),
$$
where \(\mathcal{F}_x=\mathrm{span}\{1,x,\reflenx^2,\reflenx^3\}\) and \(\mathcal{V}_r\) is the space of \(\bm{r}\)-dependent vector fields on the section. 
%
The minimum number of non-rigid warping shapes is then the \emph{separation rank} (CP rank) of \(\mathbf{R}\):
$$
\ShapeRank \;=\; \operatorname{rank}_\text{sep}\big(\mathbf{R}\big),
$$
i.e., the smallest \(N\) so that \(\mathbf{R}(\reflenx,\bm{r})=\sum_{j}^N \alpha_j(\reflenx)\,\bm{\Phi}_j(\bm{r})\) with each \(\bm{\Phi}_j\) independent of \(\reflenx\).


After assigning to $\bm{x}({x})$ and $\bm{\theta}({x})$ all rigid patterns (\(\reflenx\)-constants in $\{\bm{r},\,\FrameBasis\otimes\bm{r},\, \FrameBasis\}$, and $\FrameBasis \times \bm{r}$ ), the Saint~Venant field leaves exactly following \(\reflenx\)-independent \(\bm{r}\)-shapes:
\[
\mathcal{S}=\Big\{
\ \bm{W}_{(\varepsilon)},
\ \bm{W}_{(a)},
\ \bm{W}_{(b)},
\ \FrameBasis\otimes\bm{\WarpShearIesan},
\ \FrameBasis\;\varphi
\ \Big\},
\]
where
\begin{subequations}
\begin{alignat}{8}
&\text{Transverse:}\notag\\
&&&\bm{W}_{(\varepsilon)}:=-\,\nu\,\bm{r}&&\text{ Poisson extension }& \text{1 amplitude }&&\alpha_{\varepsilon}\propto a_\varepsilon \sim \reflenx^0\\
&&&\bm{W}_{(a)}:=-\,\nu\,\operatorname{dev}\bm{r}\otimes\bm{r}&&\text{ Poisson bending }&\text{2 amplitudes }&&\bm{\alpha}_a\propto\IesanLoadM\sim \reflenx^0 \label{eq:def-Wa}\\
&&&\bm{W}_{(b)}:=-\,\nu\big(\operatorname{dev}\bm{r}\otimes\bm{r}-\bm{r}\otimes\CentroidLocation\big)&&\text{ Poisson flexure }&\text{2 amplitudes }&&\bm{\alpha}_{b}\propto\bm{b}_{\Section}\sim \reflenx^1
\label{eq:def-Wb}
\\
%
&\text{Longitudinal:} \notag\\
&&&\FrameBasis\otimes\bm{\WarpShearIesan}&&\text{ Flexure }&\text{2 amplitudes }&&\bm{\alpha}_{\WarpShearIesan}\propto\bm{b}_{\Section}\sim \reflenx^0 \\
&&&\FrameBasis\,\WarpTwistIesan &&\text{ Torsion }&\text{1 amplitude }&&\alpha_{\WarpTwistIesan} \propto a_{\vartheta}+c_{\vartheta}\sim \reflenx^0 \\
\end{alignat}
\label{eq:def-poisson-warping}
\end{subequations}

That is \textbf{8} linearly independent shapes in the quotient \(\mathcal{V}_r/\mathscr{R}(\Section)\) for a generic section.

So, generically,
$$
\ShapeRank=8\quad(\text{or }9\text{ if you forbid absorbing the axial mode into }\bm u).
$$

\textbf{Remarks}

\begin{itemize}
\item The ``axial mode" \(+\tfrac{1}{2} \reflenx^2(\FrameBasis\otimes\CentroidLocation)\,\bm{b}_{\Section}\) is a {rigid} pattern (uniform over the section), hence it belongs to \(\mathscr{R}(\Section)\) and can be absorbed into \(\bm u(\reflenx)\). 
If, for other reasons, you choose not to absorb it (e.g., to align \(\epsilon\) with a specific constitutive form), then you would add one more shape \(\mathbf{1},\FrameBasis\) and the bound becomes \(9\). 
From a purely mathematical minimality standpoint, it does not need to be a warping shape.

\item The two Neumann–Poisson solutions \(\bm\psi\) are linearly independent modulo \(\{1,y,z\}\). 
No linear combination of \(\{1,y,z\}\) reproduces a nontrivial \(\Psi\) that satisfies the given PDE/BCs.

\item \(\varphi\) is harmonic with a different Neumann datum; it is independent from the \(\bm\WarpShearIesan\)-pair and from \(\{1,y,z\}\).

\item The three transverse Poisson blocks are quadratic (or linear) in \(\bm{r}\) and live in the \(\Section\)-plane; none is in the rigid span \(\{\PlaneBasis_y,\PlaneBasis_z,\FrameBasis\times\bm{r}\}\), and the ``(a)" and ``(b)" blocks differ by the \(\bm{r}\otimes\CentroidLocation\) term, so they are independent in non-centroidal coordinates \(\CentroidLocation \ne \mathbf{0}\).

\item Independence persists under generic boundary geometry; special symmetries (e.g., circular or some doubly symmetric sections) can collapse the rank (then \(\ShapeRank\) can be smaller).

\item Special cases that reduce \(\ShapeRank\)
\begin{itemize}
\item \textbf{Symmetric sections:} certain Poisson blocks can vanish or become dependent.
\item \textbf{Special Coordinates} \(\CentroidLocation=\mathbf{0}\) for centroidal axes removes the \(\bm{r}\otimes\CentroidLocation\) addendum.
\item \textbf{Saint–Venant–Kirchhoff circle:} \(\WarpShearIesan\)-pair simplifies, \(\tilde{\varphi}^{(a)}\) vanishes; in extreme cases some spans collapse.
\item \textbf{Pure load subsets:} if some parameters are fixed to zero (e.g., no extension), drop the corresponding shapes.
\end{itemize}
\end{itemize}
%
Formally, you can prove independence by testing against a separating family of functionals (e.g., take cross-section integrals of each component against \(\{1,y,z,y^2, yz, z^2\}\) and use integration by parts with the known PDEs/BCs of \(\varphi,\WarpShearIesan\)) to show no nontrivial combination of the eight shapes lies in \(\mathscr{R}(\Section)\).


\paragraph{Field Rank}

Here it is shown that, even after enslaving amplitudes to the alignment strain measures \(\bm{\gamma},\bm{\kappa}\) whenever their
\(\reflenx\)-dependence coincides exactly with a component of \(\bm{\gamma}\) or \(\bm{\kappa}\), at least five warping amplitudes remain independent.

The minimal count above ignores any extra constraints tying warping amplitudes to \(\bm\gamma,\,\bm\kappa\). 
These identifications restrict the admissible amplitude functions \(\alpha_j(\reflenx)\), and \(\FieldRank\) can only \emph{increase}.

Any further requirement to eliminate amplitudes to \(\bm\gamma\) and/or \(\bm\kappa\) is an additional constraint that generally increases the number of shapes beyond this lower bound.

\subparagraph{Setting}
Write the displacement as \(\hat{\bm x}=\bm x+\bm\theta\times\bm{r}+\bm\Phi(\bm{r})\,\bm\alpha(\reflenx)\), with \(\bm\Phi\) collecting
\(\reflenx\)-independent shapes. 
From \cref{prop:min8}, after removing the rigid patterns
\(\mathcal{R}(\Section)\) in \eqref{eq:rigid-span},
the Saint–Venant residual requires, for a generic section, the eight standard shapes \eqref{eq:def-poisson-warping}.

\begin{proposition}[Field rank]
For a generic cross–section (no symmetry assumptions, arbitrary \(\CentroidLocation\)), the full Saint–Venant family
\((\IesanLoadM,a_\varepsilon,a_\vartheta,\bm{b}_{\Section})\) can be represented with exactly \(\FieldRank = 5\) independent warping amplitude functions of \(\reflenx\)  once the deformations \(\bm{\gamma}(\reflenx),\bm{\kappa}(\reflenx)\) are used to associate as many warping amplitudes as possible. 
\end{proposition}

\subparagraph{Proof}
Consider the six Ie\c{s}an parameters \((a_y,a_z,a_\varepsilon,a_\vartheta,b_y,b_z)\) and their induced
\(\reflenx\)-laws in the Saint–Venant field. 
The following three blocks of warping content cannot be eliminated by \(\bm{\gamma}\) or \(\bm{\kappa}\)
without violating exactness for arbitrary parameters:

\medskip
\emph{(i) The \(a\)-Poisson in–plane pair \(\bm{W}^a_y,\bm{W}^a_z\) (two amplitude fields).}
Their amplitudes must be \((a_y,a_z)\), each \emph{constant} in \(\reflenx\).
For arbitrary \(\bm{b}_{\Section}\neq \mathbf{0}\), every component of \(\bm{\kappa}_\perp(\reflenx)\) has the form
\(C(a)+x\,D(b)\) with \(D(b)\not\equiv 0\);
every component of \(\bm{\gamma}_\perp(\reflenx)\) depends only on \(\bm{b}_{\Section}\) (not on \(\IesanLoadM\)).
Hence there is no scalar \(c\) such that \(c\,\kappa_\perp(\reflenx)\) or \(c\,\bm{\gamma}_\perp(\reflenx)\) equals a nonzero constant in \(\reflenx\) for all
\((\IesanLoadM,\bm{b}_{\Section})\). Therefore the two amplitudes tied to \(\bm{W}^{(a)}\) are \emph{independent}.

\medskip
\emph{(ii) The \(b\)-Poisson in–plane pair \(\bm{W}^b_y,\bm{W}^b_z\) (two amplitude fields).}
Their amplitudes must be \emph{linear} in \(\reflenx\), proportional to \(\bm{b}_{\Section}\), and contain no \(a\)-contribution:
\(\alpha_{P^b}(\reflenx)=\reflenx \,\mathcal{B}\,\bm{b}_{\Section}\) for a fixed invertible \(\mathcal{B}\).
Every component of \(\bm{\kappa}_\perp(\reflenx)\) equals \(C(a)+x\,D(b)\).

\medskip
\emph{(iii) The torsional axial warping \(\varphi\FrameBasis\) (one amplitude field).}
Its amplitude is \(a_\vartheta+c_\vartheta\), whereas \(\varkappa=a_\vartheta-c_\vartheta\).
Consequently, the \(\varphi\)-amplitude cannot be associated to \(\varkappa\) or \(\epsilon\) and is \emph{independent}.

\medskip
These three blocks contribute \(2+2+1=5\) independent warping amplitude functions of \(\reflenx\). 
All remaining warping amplitudes \emph{can} be associated:
the Poisson extension amplitude equals \(\epsilon=a_\varepsilon\);
the axial linear shapes \(\FrameBasis\otimes\bm{r}\) and the \(\WarpShearIesan\)-pair can be taken proportional to
\(\bm{\gamma}_\perp(\reflenx)\) (with the same \(\reflenx\)-law), eliminating them.

\subparagraph{Conclusion}

In general
\[
\text{free‑amplitude rank} \;\le\; \text{shape rank},
\]
with strict inequality whenever multiple shapes can be driven by the same \(x\)-signal taken from (\(\epsilon,\gamma,\kappa)\).

No fewer than five independent warping amplitude functions are possible for the generic full Saint–Venant family with \(\reflenx\)-independent shapes.
Hence the lower bound \(\FieldRank \ge 5\) holds, and it is sharp.

